% GENERATED FILE. Edit canonical lessons, then run npm run generate:books.
% canonical-source-hash: fnv1a64:e8bf7f4ada3c6f56
% canonical-lessons: JA-C207-tenki, JA-C207-furu, JA-C207-yuki, JA-C207-hareru, JA-C207-kumoru

\chapter{Tomorrow's Rain and Snow}
\label{ch:ja-tomorrow-s-rain-and-snow}
\hlchaptermodality{207}

\begin{chapteropening}
\textbf{By the end of this chapter:} \emph{I can say what the weather will do tomorrow.}

Complete the last lesson of chapter 207: I can say what the weather will do tomorrow.
\end{chapteropening}

\section[\texorpdfstring{\ja{てんき} (tenki)}{てんき (tenki)}]{\ja{てんき} (tenki) — the weather; fine weather}

\label{lesson:JA-C207-tenki}



\begin{quote}
Before the new one: say the Japanese for a passenger, then the Japanese for a booking.
\end{quote}

\subsection*{You'll want to know: \ja{てんき}}
\textbf{\ja{てんき}} — \emph{tenki} — \textquotedblleft{}the weather; fine weather\textquotedblright{}.

\begin{quote}
\textbf{\ja{あしたの} \ja{てんきは} \ja{いいでしょう。}} — \emph{ashita no tenki wa ii deshō} — \textquotedblleft{}The weather will probably be good tomorrow.\textquotedblright{}
\end{quote}

\begin{cousinweb}
\textbf{\ja{てんき}} (\emph{tenki}) is Sino-Japanese (kango): two Chinese characters read with their on'yomi, readings borrowed from Chinese. They are \emph{ten}, \textquotedblleft{}heaven, sky\textquotedblright{} (Middle Chinese \emph{then}) and \emph{ki}, \textquotedblleft{}air, breath\textquotedblright{} (\emph{khj+jH}) (KANJIDIC2; Middle Chinese in Baxter's transcription).
\end{cousinweb}

\begin{grammarlens}[title={What you've built}]
One more word for everyday life at A2.
\end{grammarlens}

\subsection*{Your turn}
\begin{itemize}
\raggedright
  \item \emph{Say it:} \emph{tenki}
  \item \emph{Say it:} \emph{tenki}, once more
  \item \emph{Say it:} \emph{ashita no tenki wa ii deshō} — The weather will probably be good tomorrow.
  \item \emph{Recall:} say the Japanese for a passenger, then the Japanese for a booking, then say \emph{tenki} again
\end{itemize}

\begin{tcolorbox}[breakable,colback=teal!4,colframe=teal!35!black,title={Before you move on}]
What does \ja{てんき} mean? (\textquotedblleft{}The weather; fine weather\textquotedblright{}.) Say it once more.
\end{tcolorbox}
\section[\texorpdfstring{\ja{ふる} (furu)}{ふる (furu)}]{\ja{ふる} (furu) — to fall (of rain or snow)}

\label{lesson:JA-C207-furu}



\begin{quote}
Before the new one: say the Japanese for a booking, then the Japanese for the weather.
\end{quote}

\subsection*{You'll want to know: \ja{ふる}}
\textbf{\ja{ふる}} — \emph{furu} — \textquotedblleft{}to fall (of rain or snow)\textquotedblright{}.

\begin{quote}
\textbf{\ja{あした} \ja{あめが} \ja{ふります。}} — \emph{ashita ame ga furimasu} — \textquotedblleft{}It will rain tomorrow.\textquotedblright{}
\end{quote}

Polite \textbf{\ja{ふります}}, \emph{furimasu}; a godan verb.

\begin{cousinweb}
\textbf{\ja{ふる}} (\emph{furu}) is a native Japanese word (wago), inherited rather than borrowed from Chinese; its kanji was attached later as a kun'yomi, and its older origin is not settled (JMdict, KANJIDIC2).
\end{cousinweb}

\begin{grammarlens}[title={What you've built}]
One more word for everyday life at A2.
\end{grammarlens}

\subsection*{Your turn}
\begin{itemize}
\raggedright
  \item \emph{Say it:} \emph{furu}
  \item \emph{Say it:} \emph{furu}, once more
  \item \emph{Say it:} \emph{ashita ame ga furimasu} — It will rain tomorrow.
  \item \emph{Recall:} say the Japanese for a booking, then the Japanese for the weather, then say \emph{furu} again
\end{itemize}

\begin{tcolorbox}[breakable,colback=teal!4,colframe=teal!35!black,title={Before you move on}]
What does \ja{ふる} mean? (\textquotedblleft{}To fall (of rain or snow)\textquotedblright{}.) Say it once more.
\end{tcolorbox}
\section[\texorpdfstring{\ja{ゆき} (yuki)}{ゆき (yuki)}]{\ja{ゆき} (yuki) — snow}

\label{lesson:JA-C207-yuki}



\begin{quote}
Before the new one: say the Japanese for the weather, then the Japanese for to fall.
\end{quote}

\subsection*{You'll want to know: \ja{ゆき}}
\textbf{\ja{ゆき}} — \emph{yuki} — \textquotedblleft{}snow\textquotedblright{}.

\begin{quote}
\textbf{\ja{らいしゅう} \ja{ゆきが} \ja{ふります。}} — \emph{raishū yuki ga furimasu} — \textquotedblleft{}It will snow next week.\textquotedblright{}
\end{quote}

\begin{cousinweb}
\textbf{\ja{ゆき}} (\emph{yuki}) is a native Japanese word (wago), inherited rather than borrowed from Chinese; its kanji was attached later as a kun'yomi, and its older origin is not settled (JMdict, KANJIDIC2).
\end{cousinweb}

\begin{grammarlens}[title={What you've built}]
One more word for everyday life at A2.
\end{grammarlens}

\subsection*{Your turn}
\begin{itemize}
\raggedright
  \item \emph{Say it:} \emph{yuki}
  \item \emph{Say it:} \emph{yuki}, once more
  \item \emph{Say it:} \emph{raishū yuki ga furimasu} — It will snow next week.
  \item \emph{Recall:} say the Japanese for the weather, then the Japanese for to fall, then say \emph{yuki} again
\end{itemize}

\begin{tcolorbox}[breakable,colback=teal!4,colframe=teal!35!black,title={Before you move on}]
What does \ja{ゆき} mean? (\textquotedblleft{}Snow\textquotedblright{}.) Say it once more.
\end{tcolorbox}
\section[\texorpdfstring{\ja{はれる} (hareru)}{はれる (hareru)}]{\ja{はれる} (hareru) — to clear up; to be sunny}

\label{lesson:JA-C207-hareru}



\begin{quote}
Before the new one: say the Japanese for to fall, then the Japanese for snow.
\end{quote}

\subsection*{You'll want to know: \ja{はれる}}
\textbf{\ja{はれる}} — \emph{hareru} — \textquotedblleft{}to clear up; to be sunny\textquotedblright{}.

\begin{quote}
\textbf{\ja{あしたの} \ja{あさは} \ja{はれます。}} — \emph{ashita no asa wa haremasu} — \textquotedblleft{}It will be sunny tomorrow morning.\textquotedblright{}
\end{quote}

Polite \textbf{\ja{はれます}}, \emph{haremasu}; an ichidan verb.

\begin{cousinweb}
\textbf{\ja{はれる}} (\emph{hareru}) is a native Japanese word (wago), inherited rather than borrowed from Chinese; its kanji was attached later as a kun'yomi, and its older origin is not settled (JMdict, KANJIDIC2).
\end{cousinweb}

\begin{grammarlens}[title={What you've built}]
One more word for everyday life at A2.
\end{grammarlens}

\subsection*{Your turn}
\begin{itemize}
\raggedright
  \item \emph{Say it:} \emph{hareru}
  \item \emph{Say it:} \emph{hareru}, once more
  \item \emph{Say it:} \emph{ashita no asa wa haremasu} — It will be sunny tomorrow morning.
  \item \emph{Recall:} say the Japanese for to fall, then the Japanese for snow, then say \emph{hareru} again
\end{itemize}

\begin{tcolorbox}[breakable,colback=teal!4,colframe=teal!35!black,title={Before you move on}]
What does \ja{はれる} mean? (\textquotedblleft{}To clear up; to be sunny\textquotedblright{}.) Say it once more.
\end{tcolorbox}
\section[\texorpdfstring{\ja{くもる} (kumoru)}{くもる (kumoru)}]{\ja{くもる} (kumoru) — to cloud over}

\label{lesson:JA-C207-kumoru}



\begin{quote}
Before the new one: say the Japanese for snow, then the Japanese for to clear up.
\end{quote}

\subsection*{You'll want to know: \ja{くもる}}
\textbf{\ja{くもる}} — \emph{kumoru} — \textquotedblleft{}to cloud over\textquotedblright{}.

\begin{quote}
\textbf{\ja{あしたの} \ja{ごごは} \ja{くもります。}} — \emph{ashita no gogo wa kumorimasu} — \textquotedblleft{}It will cloud over tomorrow afternoon.\textquotedblright{}
\end{quote}

Polite \textbf{\ja{くもります}}, \emph{kumorimasu}; a godan verb.

\begin{cousinweb}
\textbf{\ja{くもる}} (\emph{kumoru}) is a native Japanese word (wago), inherited rather than borrowed from Chinese; its kanji was attached later as a kun'yomi, and its older origin is not settled (JMdict, KANJIDIC2).
\end{cousinweb}

\begin{grammarlens}[title={What you've built}]
One more word for everyday life at A2.
\end{grammarlens}

\subsection*{Your turn}
\begin{itemize}
\raggedright
  \item \emph{Say it:} \emph{kumoru}
  \item \emph{Say it:} \emph{kumoru}, once more
  \item \emph{Say it:} \emph{ashita no gogo wa kumorimasu} — It will cloud over tomorrow afternoon.
  \item \emph{Recall:} say the Japanese for snow, then the Japanese for to clear up, then say \emph{kumoru} again
\end{itemize}

\begin{tcolorbox}[breakable,colback=teal!4,colframe=teal!35!black,title={Before you move on}]
What does \ja{くもる} mean? (\textquotedblleft{}To cloud over\textquotedblright{}.) Say it once more.
\end{tcolorbox}
